\section{Checks of the integrated amplitudes}
The following checks test the analytic integrations derived in Chapter 3.
Numerical values do not fix their constants or physical branches.

\subsection{Checks of the spanning matrix}
Numerical integration of the earlier two-parameter formula provides
a check of the analytic box, not its definition. At
$(s,c)=(1.6,0.2),(10,-0.4),(100,0.7)$, comparison of all
four physical polarization matrices in both photon orderings with
128-point integration in each parameter gives maximal relative
matrix differences $1.9\times10^{-13}$,
$7.1\times10^{-13}$ and $5.0\times10^{-12}$, respectively.
Separately, evaluating the rational interference coefficients
agrees with explicit covariant polarization contractions of the
closed $4\times4$ matrices to relative precision better than
$7.1\times10^{-15}$ at these points. Neither comparison uses a
published total-cross-section formula.

\subsection{Evaluation of the spanning angular primitives}
The numerical evaluation of the GPL expression can be performed
by convergent Taylor expansions without integrating the original
angle-dependent correction. Around a point $v_0$ away from a
letter $a$, write the Taylor coefficients of the shorter word
as $t_n$ and those of the derivative of $G(a,w)$ as $d_n$.
Its differential equation gives the recursion
\begin{equation}
 d_{-1}=0,\qquad
 d_n=\frac{t_n-d_{n-1}}{v_0-a},\qquad
 [ (v-v_0)^{n+1}]G(a,w;v)=\frac{d_n}{n+1}.
\end{equation}
Overlapping disks along the upper path then continue the same
analytic function to the endpoint. This operation evaluates the closed special functions on the
branch already fixed by the Feynman pole displacement and the
continuation path derived above. No cross-section comparison
enters that choice.

As a check, take $\lambda=10^{-5}$ and define
$J_j=\mathcal F_j(v_{\max};x)-\mathcal L_j(x)$ for $j=s,u$.
This is the quantity before multiplication by $1/(2x^2)$. Evaluation of the closed expressions gives
\begin{align}
 J_s(0.6)&=-18.6105396630505172566260787241204817816433900769,\\
 J_u(0.6)&=-14.1391278522157817227769489975337484383516616217,\\
 J_u(9)&=-3080.09775949773657530083054013237461366392218650.
\end{align}
Independent quadrature of the already derived angular interference
agrees with these values to absolute differences
$9\times10^{-47}$, $10^{-62}$ and $2\times10^{-58}$,
respectively. The last GPL evaluation crosses the apparent
$\eta$ pole by the specified upper path. Its residual imaginary
part is below $3\times10^{-64}$ before taking any real part.
These checks test the analytic integrations and the causal branch
choice. They do not use a published total cross section.

\subsection{Checks of the finite soft function}
The script \texttt{compton\_soft\_analytic.py} evaluates
\eqref{eq:comp-soft-K-dilog} directly and separately integrates
\eqref{eq:comp-soft-I-definition} for a check. At 70 decimal digits,
checks at $s=1.6$, $10$, $100$, and $19570.511835591835$, with
$a^2/q^2=0.0001$, $0.2$, $0.8$, and $0.9999$, give absolute differences below
$1.2\times10^{-66}$. These quadratures check the derived closed expression.
They do not define it. The result completes the finite soft parameter
integral at fixed angle. Its Born-weighted elastic angular integral
is a further part of the total-cross-section calculation.

\subsection{Check of the integrated infrared coefficient}
Direct checks of the
original angular integral at 65-digit working precision give
$\mathcal C_{\mathrm{IR}}=0.2034909958483591$, $10.14476601141435$,
$64.23179863499513$, and $359.2271977981074$ at $s=1.6$, $10$, $100$, and $19570.511835591835$. The largest absolute difference between the dilogarithmic
expression and quadrature is $1.1\times10^{-61}$.

\subsection{Independent numerical angular integration}
The numerical angular check evaluates the covariant Dirac trace before
integrating over the scattering angle.  Its scalar one-loop functions are
the logarithms and dilogarithms obtained from the parameter integrals.
It does not evaluate an integrated GPL primitive inside the numerical
integrand. In units $m=1$, put $x=s-1$, $y=x+t$, and use the
spin-averaged interference trace $\mathcal T$ and the Born factor
$\mathcal B=\overline{|\mathcal M_0|^2}/e_{\mathrm T}^4$ from
\eqref{eq:comp-supp-independent-total}.
The finite virtual-plus-soft coefficient in the local real-emission
subtraction scheme is
\begin{equation}
 \mathcal C_{\mathrm{VS}}(s)=\frac{1}{s}\int_{-1}^{1}\dd c\,
 \left\{\frac{\operatorname{Re}\mathcal T(s,c;\lambda)}4
 +\mathcal B(s,c)\left[(b_e-1)\ln\frac2\lambda-K(s,c)\right]\right\}.
\end{equation}
The logarithm of the photon mass cancels inside this combined integrand.
This coefficient must be combined with the real contribution defined by
the same local subtraction.  The finite coefficient obtained with the
analytic factor $(xy)^\eta$ has a different soft matching term. Here
$x,y$ in the regulator are the real-emission Dalitz variables, not the
elastic invariants used in this angular check.

To resolve both angular endpoints at large $s$, write $L=\ln s$ and
introduce $0<u<1$ by
\begin{align}
 c&=\left(1+\frac2x\right)\tanh\bigl[L(u-\tfrac12)\bigr],
 &y&=x-\frac{x^2}{2s}(1-c),\\
 \frac{dc}{du}&=\left(1+\frac2x\right)
 \frac{L}{\cosh^2[L(u-\tfrac12)]}.
\end{align}
The identity $\tanh(L/2)=x/(x+2)$ proves that the endpoints map exactly
to $c=-1$ and $c=1$.  Positive Gauss--Legendre nodes lie strictly inside
this interval, so no removable endpoint expression is evaluated directly.
The check uses arbitrary-precision arithmetic and three successively
doubled node counts.  Ten times the larger of the final two differences
is an empirical discretization error estimate, not a rigorous error bound.
The reported uncertainty also has the conservative arithmetic floor
$10^{20-d}$ when using $d$ decimal digits.
Increasing the arithmetic precision provides a separate check of
cancellation error, particularly near threshold.  A second variable map,
$y=(x/s)\exp(u\ln s)$, checks the angular Jacobian independently.
The scope of this test is the angular integration and its implementation.
The earlier comparison with parameter quadrature tests the scalar
one-loop functions themselves.

\subsection{Checks of the soft angular primitive}
The algebraic reduction has also been implemented independently of the
quadrature check. At $s=10$ it generates 306 terms involving 305 distinct
nonempty polylogarithm words, including words needed to evaluate their
lower-weight derivatives. Partial-fraction reconstruction of the
rational kernels agrees to $1.3\times10^{-60}$ at 60-digit precision.
At $z=0.9Z$, evaluating the generated polylogarithms through their defining
differential equations gives $17.00560385716334$, compared with
$17.00560385716422$ from direct quadrature of the rationalized kernel.
These two numerical evaluations check the algebraically integrated expression.

\subsection{Series evaluation of the soft angular endpoint}
\paragraph{Arbitrary-precision polylogarithm evaluation.}
The generated expression can be checked beyond machine precision by
convergent power series, without integrating the Compton kernel. For a
word $(a,\boldsymbol b)$ expand $G(a,\boldsymbol b;z_0+h)$ in powers of
$h$ with coefficients $c_n$, and write the tail word as
$G(\boldsymbol b;z_0+h)=\sum_n d_nh^n$. Its differential equation gives
\begin{equation}
 c_{n+1}=\frac{d_n-nc_n}{(z_0-a)(n+1)},\qquad z_0\ne a.
\end{equation}
Words are processed in increasing length, so all $d_n$ are already known.
At the regular origin of a word beginning with zero, $d_0=0$ and
$c_n=d_n/n$ for $n\geq1$. Steps smaller than the distance to the nearest
letter analytically continue these series along the real physical path.
At 60-digit working precision, the 306-term expression at $s=10$ and
$z=0.9Z$ agrees with direct quadrature to $1.0\times10^{-59}$.
Thus the finite polylogarithmic reduction itself, rather than only the
preceding integral transformations, has been checked independently.
For the total angular coefficient, taking
$z/Z=1-10^{-25}$ gives $d=8.8303688\times10^{-26}$.
Adding the analytically derived endpoint tail then yields
$\mathcal C_K(10)=10.058462232631976837745\ldots$.
Its difference from the original integral is $7.8\times10^{-48}$,
consistent with the $O(d^2\ln d)$ tail remainder. At $s=1.6$ the
interior endpoint check agrees to $2.8\times10^{-58}$.
These checks use only the amplitudes, soft integral, and algebraic
reductions derived here, with no published total cross section as input.

\subsection{Check of the inner-energy Laurent coefficients}
The script \nolinkurl{compton_real_uv_endpoint.py} verifies the
Laurent coefficients against the exact beta--hypergeometric
expressions at 80 decimal digits. This check evaluates analytic
functions and does not use numerical quadrature.

\subsection{Checks of the UV difference}
At $H=9$, the generated endpoint expression for this pole coefficient
contains six terms, with words drawn from
$G(r_\pm-1,-1;9)$, $G(r_\pm-1;9)$, and $G(-1;9)$, together with a
rational endpoint term. Evaluating this analytic expression gives
\begin{equation}
 \int_1^{10}\mathcal P_{\mathrm{UV}}(9,v)\,\dd v
 =27551.4550049022114364696640244483\ldots.
\end{equation}
The difference from independent quadrature is $4.4\times10^{-61}$ at
65-digit working precision. The quadrature uses the analytic local series
near $v=1$ to avoid subtracting individually divergent numerical terms.
The integrated result includes the complete difference with its Taylor
subtraction, rather than either divergent component in isolation.

The finite coefficient is reduced by the same rules to 26 terms.
Its origin expansions at two independent values of $H$ are
\begin{align}
 \mathcal F_{\mathrm{UV}}(9,1+\zeta)
 &=\left(\frac{24064}{729}-\frac{28672\ln3}{243}\right)\zeta+O(\zeta^2),\\
 \mathcal F_{\mathrm{UV}}(3/5,1+\zeta)
 &=\left(\frac{34816}{15}+\frac{114688}{135}\ln\frac53\right)\zeta+O(\zeta^2).
\end{align}
The integrated endpoint values are
\begin{align}
 \int_1^{10}\mathcal F_{\mathrm{UV}}(9,v)\,\dd v
 &=114586.9312643746457220611778438924\ldots,\\
 \int_1^{8/5}\mathcal F_{\mathrm{UV}}(3/5,v)\,\dd v
 &=504.4462882483088914154537366764\ldots.
\end{align}
At 65-digit precision, independent quadrature differs by
$6.3\times10^{-60}$ and $4.2\times10^{-62}$, respectively.
These values concern this angular difference sector before the overall
normalization and the other sectors are assembled.
The callable routine \texttt{uv\_difference\_closed} returns its pole or
finite coefficient from the generated polylogarithmic expression.
It performs no quadrature.

\subsection{Check of the UW pole integral}
At $H=9$ its integrand begins as $-1152t+O(t^2)$, and
\begin{equation}
 \int_0^1g_1(9,t)\,\dd t
 =-15.9227404613352952442617815523\ldots.
\end{equation}
At 60-digit working precision this expression agrees with independent
quadrature to $2.6\times10^{-56}$.

\subsection{Checks of the UW function basis}
The eight functions in the inner finite result have also been checked
against their real logarithm and dilogarithm definitions. At $H=9$ and
$t=0.1,0.5,0.9$, their polylogarithmic translations agree to better than
$2.4\times10^{-45}$ at 45-digit working precision. This includes the
difference of two dilogarithms with large negative arguments near $t=0$.
The identity is fixed by differentiation and its value at the origin,
so the numerical check does not choose a physical branch.

\subsection{Checks of the finite UW integral}
Its endpoint values before the common overall normalization are
\begin{align}
 \mathcal F_{UW}(9)
 &=-599.2785090120818758407150390268563\ldots,\\
 \mathcal F_{UW}(3/5)
 &=5.33372142619549103751931327587136\ldots.
\end{align}
Independent coordinate quadrature differs by $2.4\times10^{-74}$ and
$1.3\times10^{-65}$, respectively. These differences are consistent with
the estimated analytic tail remainder. The production routine
\texttt{real\_uw\_closed} returns the pole and finite coefficients using
only the polylogarithmic expression and convergent endpoint series.
It does not include the crossed angular sector.

\subsection{Check of the UW inner finite coefficient}
The independent numerical check integrates
the original $a$ expression only after subtracting its two endpoint
poles. It agrees with \eqref{eq:comp-real-UW-inner-closed} at an
interior point $H=0.6,t=0.3$ to better than $10^{-43}$ relatively.
The numerical integral is a check, not a definition of the answer.

\subsection{Checks of the soft-boundary conversion}
For the kernels $K=x^{-2}$ and $K=x^{-3}$ of
Section~\ref{sec:compton-soft-matching-check}, the comparison is as follows.
At $H=9$,
evaluation of the exact Mellin formulas at $\eta=\pm10^{-22}$ with
90-digit arithmetic reproduces these finite terms to more than 42 digits.
The exact cutoff formulas at $\zeta=10^{-30}$ agree with their displayed
asymptotic finite parts to more than 30 digits.  These evaluations test
analytic formulas and do not replace the integration by quadrature.

\subsection{Checks of the Bose subtraction correction}
The moments $P,A,B,C$ and finite parts $F_z,F_\eta$ use the definitions
of Section~\ref{sec:compton-soft-matching-check}. Here $H=s-1$,
and $D$ is the rational correction in
\eqref{eq:comp-real-bose-restoration}.
An independent evaluation at $H=9$ gives
\begin{align}
 P&=37497.9773085349211509285356121\ldots,\\
 A&=3813.01541537757944985908409943\ldots,\\
 B&=-1194.39517843100551019039928260\ldots,\\
 C&=875.197982980881880082536498852\ldots,\\
 F_\eta-F_z&=8697.24837987961462163725697550\ldots.
\end{align}
These values are evaluations of the closed moments, with no numerical
integration.  They provide the finite matching contribution of this
individual subtraction correction.

At $H=9$, the resulting closed endpoint value is
\begin{align}
 F_z[D]&=-34968.0992648172707249254370242076502947231\ldots,\\
 F_\eta[D]&=-26270.8508849376561032881800487118253811979\ldots.
\end{align}
Evaluations requested at 45 and 70 decimal digits agree in all the displayed
digits.  As an independent check only, quadrature of the one-dimensional
kernel up to $u=s-0.02$ agrees with its generalized-polylogarithm primitive
to an absolute difference below $1.1\times10^{-79}$ at 85-digit precision.
The finite endpoint value itself uses the analytic asymptotic expansion,
not quadrature or a fit to a cutoff sequence.

\subsection{Finite parts of the real-energy primitives}
Exact symbolic differentiation verifies the displayed GPL
primitives for all three sectors. As independent checks, at $H=9$
their $z$-cut finite parts are
\begin{align}
 F_{z,W}&=-199.1245079022571280440396316233659,\\
 F_{z,\ln s}&=371.2012216070278610540821023119374,\\
 F_{z,\mathrm{reg}}&=-13063.03114094500957535028266440329.
\end{align}
These numbers do not include the common factor $1/(8sH)$ or the
conversion from the $z$ cut to the symmetric analytic regulator.
Separate real integrations check the interior primitive boundary
differences to better than fifty decimal places at this point.
Continuing the full regular-sector GPL combination above and below
its artificial interior letter agrees to eighty decimal places.
These checks concern evaluation of the analytic primitives and do
not supply coefficients to them.

\subsection{Checks of the real soft residues}
The following checks evaluate the analytic residue moments and their
integrated sum. At $H=0.6,9,99$, separate 70-digit quadratures of these residue
moments agree with the analytic expressions to better than
$10^{-60}$ relatively.

At $H=9$ this correction is
\begin{equation}
 A_D=3813.01541537757944985908409942886783767351749\ldots.
\end{equation}
Evaluating the analytic expressions on both sides of
\eqref{eq:comp-real-integrated-pole-check} gives
\begin{center}
\begin{tabular}{rll}
\toprule
$H$ & Common value & Relative difference\\
\midrule
$0.6$ & $0.976756780072123721541847283885$ & $7.0\times10^{-68}$\\
$9$ & $730.423152821833269374946708743$ & $5.0\times10^{-70}$\\
$99$ & $50871.5845189161404834801869055$ & $6.0\times10^{-66}$\\
\bottomrule
\end{tabular}
\end{center}
This is a numerical check of the separately derived analytic
coefficients at three energies. It does not replace an identity
proof for arbitrary $H$, nor does it verify the finite total.

\subsection{Fixed-energy checks of endpoint subtraction}
At $H=9$ the pole integrand has the local expansion
\begin{equation}
 \mathcal P_{\mathrm{UV}}(9,1+\zeta)
 =-\frac{7168}{243}\zeta+O(\zeta^2).
\end{equation}
Thus its apparent origin poles cancel and the lower endpoint is regular.

As an independent check, the coefficient of its analytic-regulator pole
at $H=9$ integrates to
\begin{equation}
 P_{\mathrm{shift}}(9)
 =20\ln10+\frac{2913280}{59319}\ln\frac{11}{2}
 +\frac{5602208}{1521}.
\end{equation}
This term checks the subtraction bookkeeping independently of the
finite-angle polylogarithm evaluation.
